\section{Ethics and Social Reasoning in AI}
\label{sec:9.1}
Chapter~\ref{sec:2} laid out the ethical frameworks a design might draw on, how a principle gets translated into an objective function, and the neuroscience of empathy and theory of mind; chapter~\ref{sec:4} covered how a system might learn from experience and observation. Both left open problems. The first is moral responsibility: who is answerable when a system's ethical behavior fails, given that no single party fully controls it. The second asks whether the methods for embedding a principle in a system generalize past the cases they were built and tested on, or only look like they do. The third covers value learning and alignment as active technical research rather than settled method. Two more return to the neuroscience of empathy and theory of mind from chapter~\ref{sec:2} and the learning and social-reasoning questions from chapters~\ref{sec:4} and \ref{sec:5}, asking what remains unknown there and what a research program building on it would need to investigate next. The sixth takes up the one this book has the most at stake in: what an antifascist AI system would actually have to detect, and why the capability chapters~\ref{sec:4} through \ref{sec:6} assume it has is further off than those chapters imply. The last takes up what chapter~\ref{sec:3} leaves owing: what is owed to the bearer a floor requires, what a guardianship regime for one would look like, and whether anyone outside the arrangement could check it.
